\begin{frame}{협력 vs 경쟁, 그리고 self-play}
  MARL 문제는 크게 \textbf{두 극단}으로 나뉨:
  \begin{columns}[T]
    \begin{column}{0.48\textwidth}
      \kb{완전 협력}
      \\[0.4em]
      {\footnotesize 공통 보상을 공유하는 여러 에이전트
      (예: 물류창고의 로봇 여러 대가 함께 짐을 옮기는
      문제).}
    \end{column}
    \begin{column}{0.48\textwidth}
      \kb{완전 경쟁}
      \\[0.4em]
      {\footnotesize 제로섬 게임 --- 한쪽이 얻으면 다른
      쪽이 정확히 그만큼 잃음(예: 바둑, 체스). 실전 문제는
      이 \textbf{둘 사이 어딘가}에 많음(자율주행차: 충돌
      회피에서는 협력, 도로 공간을 두고는 경쟁 ---
      Week 14의 단일 에이전트 로봇과 대비).}
    \end{column}
  \end{columns}
  \vskip0.7em
  \begin{block}{self-play: 세 실습과 AlphaZero를 잇는
  관념}
    AlphaZero가 \textbf{스스로와 대국하며} 강해진 것도
    멀티에이전트 학습의 \textbf{특수 형태} --- 상대가 매번
    ``지금 막 학습한 \textbf{나 자신의 복사본}''이라, 상대
    실력이 내 실력과 함께 계속 높아지는 \textbf{자연스러운
    커리큘럼}.
  \end{block}
\end{frame}
